\subsection{Complete event coverage and input margins}

This follow-up audit was performed after the C/D results were known. All
500,000 launched ledger records in each catalog were read. Eligibility was
reconstructed from the saved first-visit entry/exit geometry and exit momentum,
without retracing Geant4 steps. The resulting 499,285 C and 499,275 D eligible
rows matched every saved scoring input, including event identifiers and their
order. The 715 and 725 ineligible events were excluded by the physical event
contract before scoring, not by numerical filtering. The saved target agreed
bitwise with the independently recomputed $Y=p_{\rm in}-p_{\rm out}$.

For each eligible row, define
\begin{equation}
 m_{32}=\operatorname{fl}_{32}\!\left(
 \operatorname{fl}_{32}(p_{\rm in})-\operatorname{fl}_{32}(Y)\right).
\end{equation}
Both catalogs have zero nonfinite, zero and negative margins. None has a
positive margin at or below $10^{-3}$ GeV/$c$, or at or below 1024 downward
ULPs at $p_{\rm in}^{(32)}$; the smallest margin-to-ULP ratios are
$8.2840175\times10^6$ and $8.311172\times10^6$, respectively. The cumulative
thresholds were fixed before the ledger scan. Here one downward ULP is
$p^{(32)}-\operatorname{nextafter}(p^{(32)},-\infty)$.
Table~\ref{tab:actual_support_margins} reports the distribution after computing
the subtraction in float32 and promoting its result to float64 for linear
quantile interpolation. A float64 difference of two rounded inputs is a
distinct diagnostic, retained separately in the audit records.

\begin{table}[!htbp]
\centering\small
\setlength{\tabcolsep}{4pt}
\caption{Actual float32 support-margin distributions, in GeV/$c$. Counts
include all eligible events. Zero, negative and nonfinite counts are zero in
both catalogs. The complete finer quantiles and seven momentum-bin summaries
are retained with the aggregate audit.}
\label{tab:actual_support_margins}
\begin{tabular}{lrrrrrr}
\toprule
Catalog & Events & Minimum & 1\% & Median & 99\% & Maximum\\
\midrule
C & 499,285 & 0.295779943 & 0.306910386 & 0.649518907 & 0.992840517 & 0.999970257\\
D & 499,275 & 0.297881365 & 0.306801179 & 0.650127709 & 0.992949682 & 0.999970496\\
\bottomrule
\end{tabular}
\end{table}

All 240 registered models per catalog retained the complete eligible-event
count and event-order digest. Recomputing only their frozen centers and
standardization, without evaluating a neural flow, reproduced all 960 saved
$z$ and $b$ arrays bitwise: 479,308,800 values, with no mismatch. The smallest
$b-z$ was 751.749807, and no saved $z$ reached or exceeded $b$.
All 5280 saved model-component arrays passed their recorded byte, shape and
dtype checks, and all event log densities were finite. Together these checks
show that no eligible row was silently dropped or numerically rejected from
the completed scores. They are a ledger and preprocessing reconciliation,
not a new neural-density evaluation or a validation of the transport physics.

\subsection{Actual support log-CDFs and their mathematical scale}

Across all 480 registered model--catalog evaluations, every saved support
log-CDF is float32 negative zero: 239,654,400 entries, zero nonzero values and
zero nonfinite values. Thus all its recorded quantiles are zero. These are
repeated model--event evaluations, not that many independent events; the 20
Full rerun pairs per catalog retain their exact duplicate identities.
Table~\ref{tab:actual_support_bounds} separates the primary 80 runs from the
160 earlier registered runs.

\begin{table}[!htbp]
\centering\small
\caption{Range of recorded standardized support bounds across every eligible
event and model in each cohort. Each model has bound 8 for both spline layers.
All support log-CDF entries are negative zero; all saved $z<b$.}
\label{tab:actual_support_bounds}
\begin{tabular}{llrrr}
\toprule
Catalog & Cohort & Model--event entries & Minimum $b$ & Maximum $b$\\
\midrule
C & Primary 80 & 39,942,800 & 761.355408 & 2776.529053\\
C & Earlier 160 & 79,885,600 & 761.350159 & 2776.529053\\
D & Primary 80 & 39,942,000 & 761.354797 & 2776.541260\\
D & Earlier 160 & 79,884,000 & 761.349548 & 2776.543945\\
\bottomrule
\end{tabular}
\end{table}

The zeros have a quantitative explanation. Every saved bound exceeds 761.349,
well outside the spline interval $[-8,8]$. Each layer is exactly the identity
there, so its transformed base bound equals $b$, independently of the neural
conditioner outputs. With the standard logistic base,
\begin{equation}
 \epsilon(b)=1-F_\theta(b\mid h)=\frac{1}{1+e^b},\qquad
 \eta(b)=-\log F_\theta(b\mid h)=\log(1+e^{-b}).
\end{equation}
For $b>0$, both satisfy $0<\epsilon(b),\eta(b)<e^{-b}$. The smallest recorded
bound therefore gives $\epsilon,\eta<10^{-330}$ for every evaluated entry.
The excluded probability mass and the likelihood normalization correction
are mathematically positive, but are smaller than half the smallest positive
float64 subnormal number. A float64 re-expression would also round them to
zero; it would not reveal a representable nonzero distribution. The aggregate
audit retains the $b$ quantiles and log-domain bounds, rather than assigning
artificial nonzero probabilities.

This certificate concerns the mathematical logistic function evaluated at the
saved rounded bounds. It uses the verified identity-tail construction and
does not rerun the neural networks in float64. It explains why normalization
has no observable contribution for these models and catalogs; the term remains
in the likelihood and is not assumed negligible for other models or inputs.
